% sec8_1.tex — Section 8.1: Qualitative Response (QR) Models

In many applications in economics and other sciences, the dependent variable is
a categorical variable. For example, a person may be in the labor force or not;
a commuter chooses a particular mode of transportation; a patient may die or not;
and so on. Regression models in which the dependent variable takes on discrete
values are called \textbf{qualitative response (QR) models}. A QR model is described
by a parameterized density function $f(y_t \mid \mathbf{x}_t; \boldsymbol{\theta})$ representing a family of \emph{discrete}
probability distributions of $y_t$ given $\mathbf{x}_t$, indexed by $\boldsymbol{\theta}$.

A QR model is called a \textbf{binary response model} if the dependent variable can
take on only two values (which can be taken to be 0 and 1 without loss of generality)
and is called a \textbf{multinomial response model} if the dependent variable can
take on more than two values. Only binary models are treated in this subsection.
For a thorough treatment of binary and multinomial models, see Amemiya (1985,
Chapter~9).

The most popular binary model is the probit model, which has already
been presented in Example~7.3. Another popular binary model is the \textbf{logit model}:
\begin{equation}
  \begin{cases}
    f(y_t = 1 \mid \mathbf{x}_t; \boldsymbol{\theta}_0) = \Lambda(\mathbf{x}_t'\boldsymbol{\theta}_0), \\
    f(y_t = 0 \mid \mathbf{x}_t; \boldsymbol{\theta}_0) = 1 - \Lambda(\mathbf{x}_t'\boldsymbol{\theta}_0),
  \end{cases}
  \label{eq:8.1.1} \tag{8.1.1}
\end{equation}
where $\Lambda$ is the cumulative density function of the logistic distribution:
\begin{equation}
  \Lambda(v) \equiv \frac{\exp(v)}{1 + \exp(v)}.
  \label{eq:8.1.2} \tag{8.1.2}
\end{equation}
This density $f(y_t \mid \mathbf{x}_t; \boldsymbol{\theta}_0)$ can be written as
\begin{equation}
  f(y_t \mid \mathbf{x}_t; \boldsymbol{\theta}_0)
  = \Lambda(\mathbf{x}_t'\boldsymbol{\theta}_0)^{y_t}
    \times [1 - \Lambda(\mathbf{x}_t'\boldsymbol{\theta}_0)]^{1-y_t}.
  \label{eq:8.1.3} \tag{8.1.3}
\end{equation}
Taking logs of both sides and replacing the true parameter value $\boldsymbol{\theta}_0$ by its
hypothetical value $\boldsymbol{\theta}$, we obtain the log likelihood for observation~$t$:
\begin{equation}
  \log f(y_t \mid \mathbf{x}_t; \boldsymbol{\theta})
  = y_t \log \Lambda(\mathbf{x}_t'\boldsymbol{\theta})
    + (1 - y_t) \log[1 - \Lambda(\mathbf{x}_t'\boldsymbol{\theta})].
  \label{eq:8.1.4} \tag{8.1.4}
\end{equation}

The logit objective function $Q_n(\boldsymbol{\theta})$ is $1/n$ times the log likelihood of the sample
$(y_1, \mathbf{x}_1, y_2, \mathbf{x}_2, \ldots, y_n, \mathbf{x}_n)$. Under the assumption that
$\{y_t, \mathbf{x}_t\}$ is i.i.d., the log likelihood of the sample is the sum of the log likelihood
for observation~$t$ over~$t$. Therefore,
\begin{equation}
  Q_n(\boldsymbol{\theta})
  = \frac{1}{n} \sum_{t=1}^{n}
    \bigl\{ y_t \log \Lambda(\mathbf{x}_t'\boldsymbol{\theta})
    + (1 - y_t) \log[1 - \Lambda(\mathbf{x}_t'\boldsymbol{\theta})] \bigr\}.
  \label{eq:8.1.5} \tag{8.1.5}
\end{equation}
For probit, we have proved that its log likelihood function is concave (Example~7.6),
and that the ML estimator is consistent (Example~7.9) and asymptotically
normal (Example~7.11) under the assumption that $\E(\mathbf{x}_t\mathbf{x}_t')$ is nonsingular. You
are about to find out that doing the same for logit is easier. (On the first reading,
however, you may wish to skip the discussion on consistency and asymptotic
normality.)

\subsection*{Score and Hessian for Observation \texorpdfstring{$t$}{t}}

The logistic cumulative density function has the convenient property that
\begin{equation}
  \Lambda'(v) = \Lambda(v)(1 - \Lambda(v)), \quad
  \Lambda''(v) = [1 - 2\Lambda(v)]\Lambda(v)[1 - \Lambda(v)].
  \label{eq:8.1.6} \tag{8.1.6}
\end{equation}
Using this, it is easy to derive (see Review Question~1(a)) the following expressions
for the score and the Hessian for observation~$t$:
\begin{align}
  \mathbf{s}(\mathbf{w}_t; \boldsymbol{\theta})
  \;\biggl(\equiv \frac{\partial \log f(y_t \mid \mathbf{x}_t; \boldsymbol{\theta})}
    {\partial \boldsymbol{\theta}}\biggr)
  &= [y_t - \Lambda(\mathbf{x}_t'\boldsymbol{\theta})]\,\mathbf{x}_t,
  \label{eq:8.1.7} \tag{8.1.7} \\[6pt]
  \mathbf{H}(\mathbf{w}_t; \boldsymbol{\theta})
  \;\biggl(\equiv \frac{\partial \mathbf{s}(\mathbf{w}_t; \boldsymbol{\theta})}
    {\partial \boldsymbol{\theta}'}\biggr)
  &= -\Lambda(\mathbf{x}_t'\boldsymbol{\theta})
     [1 - \Lambda(\mathbf{x}_t'\boldsymbol{\theta})]\,
     \mathbf{x}_t \mathbf{x}_t',
  \label{eq:8.1.8} \tag{8.1.8}
\end{align}
where $\mathbf{w}_t = (y_t, \mathbf{x}_t')'$.

\subsection*{Consistency}

Since $\mathbf{x}_t \mathbf{x}_t'$ is positive semidefinite, it immediately follows from its expression that
$\mathbf{H}(\mathbf{w}_t; \boldsymbol{\theta})$ is negative semidefinite and hence $Q_n(\boldsymbol{\theta})$ is concave.
The relevant consistency theorem, therefore, is Proposition~7.6. We show here that the conditions
of the proposition are satisfied under the nonsingularity of $\E(\mathbf{x}_t\mathbf{x}_t')$.
Condition~(a) of the proposition is the (conditional) density identification:
$f(y_t \mid \mathbf{x}_t; \boldsymbol{\theta}) \neq f(y_t \mid \mathbf{x}_t; \boldsymbol{\theta}_0)$
with positive probability for $\boldsymbol{\theta} \neq \boldsymbol{\theta}_0$.\footnote{In other words, let
$X = f(y_t \mid \mathbf{x}_t; \boldsymbol{\theta})$ and
$Y = f(y_t \mid \mathbf{x}_t; \boldsymbol{\theta}_0)$. $X$ and $Y$ are random variables because they
are functions of $\mathbf{w}_t = (y_t, \mathbf{x}_t')'$. The condition can be stated simply as
$\text{Prob}(X \neq Y) > 0$ for $\boldsymbol{\theta} \neq \boldsymbol{\theta}_0$.}
Since a logarithm is a strictly monotone transformation, this condition is equivalent to
\begin{equation}
  \log f(y_t \mid \mathbf{x}_t; \boldsymbol{\theta})
  \neq \log f(y_t \mid \mathbf{x}_t; \boldsymbol{\theta}_0)
  \text{ with positive probability for }
  \boldsymbol{\theta} \neq \boldsymbol{\theta}_0,
  \label{eq:8.1.9} \tag{8.1.9}
\end{equation}
where $\boldsymbol{\theta}_0$ is the true parameter value and the density $f$ is given in \eqref{eq:8.1.3}.
Verification of this condition is the same as in the probit example in Example~7.9 and
proceeds as follows. The discussion in Example~7.8 has established that
\begin{equation}
  \E(\mathbf{x}_t\mathbf{x}_t') \text{ nonsingular}
  \;\Rightarrow\;
  \mathbf{x}_t'\boldsymbol{\theta} \neq \mathbf{x}_t'\boldsymbol{\theta}_0
  \text{ with positive probability for }
  \boldsymbol{\theta} \neq \boldsymbol{\theta}_0.
  \label{eq:8.1.10} \tag{8.1.10}
\end{equation}
Since the logit function $\Lambda(v)$ is a strictly monotone function of~$v$, we have
$\Lambda(\mathbf{x}_t'\boldsymbol{\theta}) \neq \Lambda(\mathbf{x}_t'\boldsymbol{\theta}_0)$
when $\mathbf{x}_t'\boldsymbol{\theta} \neq \mathbf{x}_t'\boldsymbol{\theta}_0$.
Thus, condition~(a) is implied by the nonsingularity of $\E(\mathbf{x}_t\mathbf{x}_t')$.
Turning to condition~(b) (that
$\E\bigl[|\log f(y_t | \mathbf{x}_t; \boldsymbol{\theta})|\bigr] < \infty$ for
all $\boldsymbol{\theta}$), it is easy to show that
\begin{equation}
  |\log \Lambda(v)| \leq |\log \Lambda(0)| + |v|.
  \label{eq:8.1.11} \tag{8.1.11}
\end{equation}
Condition~(b) can then be verified in a way analogous to the verification of the
same condition for probit in Example~7.9. Thus, as in probit, the nonsingularity of
$\E(\mathbf{x}_t\mathbf{x}_t')$ ensures that logit ML is consistent.

\subsection*{Asymptotic Normality}

To prove asymptotic normality under the same condition, we verify the conditions
of Proposition~7.9. Condition~(1) of the proposition is satisfied if the parameter
space $\Theta$ is taken to be $\mathbb{R}^p$. Condition~(2) is obviously satisfied. Condition~(3) can be
checked easily as follows. Since $\E(y_t \mid \mathbf{x}_t) = \Lambda(\mathbf{x}_t'\boldsymbol{\theta}_0)$, we have
$\E[\mathbf{s}(\mathbf{w}_t; \boldsymbol{\theta}_0) \mid \mathbf{x}_t] = \mathbf{0}$
from \eqref{eq:8.1.7}. Hence,
$\E[\mathbf{s}(\mathbf{w}_t; \boldsymbol{\theta}_0)] = \mathbf{0}$ by the Law of Total Expectations. It is left
as Review Question~1(b) to derive the conditional information matrix equality
\begin{equation}
  \E\bigl[\mathbf{s}(\mathbf{w}_t; \boldsymbol{\theta}_0)\,
    \mathbf{s}(\mathbf{w}_t; \boldsymbol{\theta}_0)' \mid \mathbf{x}_t\bigr]
  = -\,\E\bigl[\mathbf{H}(\mathbf{w}_t; \boldsymbol{\theta}_0) \mid \mathbf{x}_t\bigr].
  \label{eq:8.1.12} \tag{8.1.12}
\end{equation}
Hence, (3) is satisfied by the Law of Total Expectations. The local dominance
condition (condition~(4)) is easy to verify for logit because, since
$|\Lambda(\mathbf{x}_t'\boldsymbol{\theta})[1 - \Lambda(\mathbf{x}_t'\boldsymbol{\theta})]| < 1$,
we have: $\|\mathbf{H}_t(\mathbf{w}_t; \boldsymbol{\theta})\| \leq \|\mathbf{x}_t\mathbf{x}_t'\|$ for all
$\boldsymbol{\theta}$. The expectation of $\|\mathbf{x}_t\mathbf{x}_t'\|^2$
(the sum of squares of the elements of $\mathbf{x}_t\mathbf{x}_t'$) is finite if
$\E(\mathbf{x}_t\mathbf{x}_t')$ is nonsingular (and hence finite).
So $\E(\|\mathbf{x}_t\mathbf{x}_t'\|)$ is finite. Regarding condition~(5), the argument in
footnote~29 of Newey and McFadden (1994) used for verifying the same condition
for probit can be used for logit as well.\footnote{Only for interested readers, here's
the argument. For any $\bar{v} > 0$, there exists a $C$ such that
$\Lambda(v)[1 - \Lambda(v)] \geq C > 0$ for $|v| \leq \bar{v}$. So
$\E\bigl[\Lambda(\mathbf{x}'\boldsymbol{\theta})[1 - \Lambda(\mathbf{x}'\boldsymbol{\theta})]\mathbf{x}\mathbf{x}'\bigr]
\geq \E\bigl[\mathbf{1}(|\mathbf{x}'\boldsymbol{\theta}| \leq \bar{v})
\Lambda(\mathbf{x}'\boldsymbol{\theta})[1 - \Lambda(\mathbf{x}'\boldsymbol{\theta})]\mathbf{x}\mathbf{x}'\bigr]
\geq C\,\E\bigl[\mathbf{1}(|\mathbf{x}'\boldsymbol{\theta}| \leq \bar{v})\mathbf{x}\mathbf{x}'\bigr]$
in the matrix inequality sense, where $\mathbf{1}(|\mathbf{x}'\boldsymbol{\theta}| \leq \bar{v})$ is the
indicator function. The last term is positive definite (not just positive semi-definite)
for large enough $\bar{v}$ by non-singularity of $\E(\mathbf{x}\mathbf{x}')$.}

Thus, we conclude that if $\{y_t, \mathbf{x}_t\}$ is i.i.d.\ and $\E(\mathbf{x}_t\mathbf{x}_t')$ is nonsingular, then the
logit ML estimator $\hat{\boldsymbol{\theta}}$ of $\boldsymbol{\theta}_0$ is consistent and asymptotically normal, with the asymptotic variance consistently estimated by
\begin{equation}
  \widehat{\avar}(\hat{\boldsymbol{\theta}})
  = -\biggl[\frac{1}{n}\sum_{t=1}^{n}
    \mathbf{H}(\mathbf{w}_t; \hat{\boldsymbol{\theta}})\biggr]^{-1},
  \label{eq:8.1.13} \tag{8.1.13}
\end{equation}
where $\mathbf{w}_t = (y_t, \mathbf{x}_t')'$ and $\mathbf{H}(\mathbf{w}_t; \boldsymbol{\theta})$ is given in \eqref{eq:8.1.8}.

\bigskip
\noindent\rule{\textwidth}{0.4pt}
\textbf{QUESTIONS FOR REVIEW}
\medskip

\begin{enumerate}
\item (Score and Hessian for observation~$t$ for general densities)\quad
  In \eqref{eq:8.1.1}, replace the logit cumulative density function (c.d.f.)
  $\Lambda(\mathbf{x}_t'\boldsymbol{\theta}_0)$ by a general c.d.f.\ $F(\mathbf{x}_t'\boldsymbol{\theta}_0)$.

  \begin{enumerate}[label=(\alph*)]
  \item Verify that the score and Hessian for observation~$t$ are given by
    \[
      \mathbf{s}(\mathbf{w}_t; \boldsymbol{\theta})
      = \frac{y_t - F_t}{F_t \cdot (1 - F_t)}\,f_t\,\mathbf{x}_t,
    \]
    \[
      \mathbf{H}(\mathbf{w}_t; \boldsymbol{\theta})
      = -\biggl[\frac{y_t - F_t}{F_t \cdot (1 - F_t)}\biggr]^2
        f_t^2\,\mathbf{x}_t\mathbf{x}_t'
        + \biggl[\frac{y_t - F_t}{F_t \cdot (1 - F_t)}\biggr]
        f_t'\,\mathbf{x}_t\mathbf{x}_t',
    \]
    where $F_t = F(\mathbf{x}_t'\boldsymbol{\theta})$,
    $f_t = f(\mathbf{x}_t'\boldsymbol{\theta})$, and
    $f(\cdot) = F'(\cdot)$ is the density function.

  \item Verify the conditional information matrix equality \eqref{eq:8.1.12} for the general
    c.d.f.\ $F$. \textbf{Hint:} Show that the expected value of the second term in the
    expression for the Hessian in the previous question is zero.
  \end{enumerate}

\item Verify \eqref{eq:8.1.7} and \eqref{eq:8.1.8}.

\item (Logit for serially correlated observations)\quad
  Suppose $\{y_t, \mathbf{x}_t\}$ is ergodic stationary but not necessarily i.i.d.
  Is the logit ML estimator described in the text consistent?
  [Answer: Yes, because the relevant consistency theorem, Proposition~7.6,
  does not require a random sample. However, the expression for
  $\avar(\hat{\boldsymbol{\theta}})$ is not given by the negative of the inverse of the sample Hessian.]
\end{enumerate}
